\documentclass{zhvt-classic}

% 使用本地源流明体（GenRyuMin2 TW）按文件路径加载，不依赖系统字体名
% 版框（内框／界栏／书口侧栏／鱼尾）改用与版心印章同一印泥棕色：
%   印章图片 yshs.png 实色 = RGB(135,68,52) = HTML 874434
\usepackage{xcolor}
\usepackage{tikz}
% sealbrown 双注册：xcolor 供 TikZ/pgf，L3 供 cls 版框
\definecolor{sealbrown}{HTML}{874434}
\ExplSyntaxOn
\color_set:nnn { sealbrown } { HTML } { 874434 }
\ExplSyntaxOff

\zhvtset{
  grid color       = sealbrown,
  main font        = GenRyuMin2TW-SB,
  main font path   = C:/Users/Administrator/texmf/fonts/opentype/public/GenRyuMin2TW/,
jiazhu font      = GenRyuMin2TW-M,
  jiazhu font path = C:/Users/Administrator/texmf/fonts/opentype/public/GenRyuMin2TW/,
  % 页码框容纳 4 字：正常页「九十七」+前置空 1 格=4 行恰好；
  % 100+ 页「一百零一」共 5 行超出，末字「一/二/三」被框裁掉（2026-10-04 实测），故放宽到 5
  shukou page number chars = 5,
  % 书口书名框高从默认 6 字减到 4 字，把省出的 2 字高拨给 chapter 框，
  % 使 6 字书名+卷名放进 9 字高的 chapter 框后不再压线（2026-10-05）
  shukou title chars = 4,
}

\title{史记}[楚世家]
\author{汉·司马迁}
\maker{闲叟漫记·重製}


% ===== 索隐述赞样式（2026-11-08 四次定稿 · 愚大人亲手改定参数）=====
% ⚠ 竖排前提：自定义字体家族必须带 RawFeature={vertical:+vert}；
%   simhei.ttf 加 +vhal 会整体加载失败静默回退，只写 +vert；simfang 同理。
\newCJKfontfamily{\shuzanfont}[Path=C:/Windows/Fonts/,RawFeature={vertical:+vert}]{simhei.ttf}
\newCJKfontfamily{\shuzanbody}[Path=C:/Windows/Fonts/,RawFeature={vertical:+vert}]{simfang.ttf}
% 题头：152x30pt 圆角 6pt 棕底白字；baseline (0,2) 使整框落入界行格子不压版框线；
% 白字节点 x 不补偿、y +8.2pt（像素级墨迹实测 dx/dy≈0，L/R 留白约 20pt）。
\newcommand{\shuzanhead}[1]{%
  \begin{tikzpicture}[baseline={(0,2)}]
    \fill[sealbrown, rounded corners=6pt]
      (-76pt,-15pt) rectangle (76pt,15pt);
    \node[text=white, inner sep=0pt, anchor=center] at (0pt,8.2pt)
      {\shuzanfont\fontsize{23}{23}\selectfont #1};
  \end{tikzpicture}%
}
% ★ 修正书口卷名位置：chapter 框起始位置改回 cls 默认值，避免全局影响。
%   只在每个卷首页通过 \juan 命令里的空命令微调该页位置（安全，不会跨页污染）。
\ExplSyntaxOn
\AtBeginDocument{
  \dim_set:Nn \l__zhvt_shukou_chapter_title_start_dim
    { \l__zhvt_shukou_title_height_dim + \l__zhvt_font_size_dim * 2 }
  \dim_set:Nn \l__zhvt_shukou_page_number_start_dim
    { \l__zhvt_shukou_chapter_title_start_dim + \l__zhvt_shukou_chapter_title_height_dim }
}
\ExplSyntaxOff

\begin{document}

% ★★ 版心（书口）排布覆写：把书名放到鱼尾之下（2026-10-04）
%
%   cls 默认版心自上而下：书名（title 静态框，在鱼尾之上）→ 鱼尾 → 篇名（chapter 动态框）→ 页码
%   本模板改为：          鱼尾（位置不动）→ 书名 → 篇名 → 页码分割线 → 页码
%
%   手法：把 cls 放在 title 静态框里的书名，改放进 chapter 动态框（鱼尾之下）；
%         title 静态框留空。书名取 \l__zhvt_title_str，即 \title{...} 的第一参数，
%         所以换书时只改 \title，不必在此重复写书名。
%
%   ⚠ 三条必须遵守的写法（2026-10-04 实测）：
%     ① 汉字必须用 \hbox{...} 包住。flowfram 会把框内容装进一个宽 14.4mm（= 一个 line_sep
%        = 书口栏宽）的纵向 minipage 并按该宽度折行，而整页靠 PDF /Rotate 90 显示、
%        框内文字本身不变向 —— 不包 \hbox 时「围炉夜话」4 字会被逐字折成 4 行，
%        在版面上就成了横排 4 字并横向溢出版框。
%     ② 竖排下**一个 \hbox 就是一整列**。书名与篇名要同处一列，必须放进同一个 \hbox；
%        分成两个 \hbox 时，后一个会被当成新的一列、甩到隔壁栏去
%        （实测「序」跑到书口栏外侧 14.4mm 的正文列上）。中间的 \quad 在框内是横向空白，
%        显示出来正好是书名与篇名之间的纵向间隔。
%     ③ ★ 内容盒要写成 **与框等宽的定宽盒** `\hbox to \l__zhvt_shukou_chapter_title_height_dim{…\hfil}`。
%        实测：chapter 框是把内容盒沿纵向**居中**摆放的，内容宽度一变，书名就跟着漂：
%        自然宽度的 \hbox（书名 38.4mm）落在 x=157.2mm，加篇名后（57.6mm）落在 x=152.4mm，
%        目录页（无篇名）与正文页的书名就差半格。做成等宽盒后，无论有没有篇名，
%        内容固定从框顶起排，实测两页均落在 x=142.8mm，不再漂。
\ExplSyntaxOn
% 把书口 chapter 框（书名+篇名）起始位置下移 1 字高，避免 4 字卷名顶到鱼尾
% ⚠ cls 的 \__zhvt_calculate_dimensions: 在 \AtBeginDocument 里执行，会重算该 dim；
%   必须把覆写挂到 begin document 之后，否则会被二次覆盖（2026-10-05 实测）。
% (旧方案已移到导言区，直接修改 chapter 框起始位置)
% 卷标题命令 #1=书口篇名(鱼尾) #2=正文卷名（含卷序，卷名内空格用 ~）
%
% ★★ 目录页两项修订（2026-10-05，愚大人指出）★★
%   ① 页码起点统一对齐：
%      原状 —— 目录项写成 \par\hspace{4em}卷名\hspace{4em}页码，页码起点 = 卷名实宽 + 4em。
%      卷名实宽随字数变：「卷一~梁惠王上」6 字 +~、「卷七~离娄上」5 字 +~，
%      差 1 个字位 → 卷七~卷十 4 条的页码整体比其余 10 条高 27.21pt（一个字），看着上下不齐。
%      改法 —— 卷名补全角空格（宽 = 1 字 = 27.215pt）补足到 **7 字符位**再写进目录，
%      14 条页码起点遂全部落在同一条竖线上（起点对齐；末字随页码字数自然参差）。
%      ⚠ 只对目录里的副本 \l_tmpc_tl 补，正文卷名仍用 \l_tmpa_tl 原样 —— 正文像素不受影响。
%   ② 卷名整体上移半字位：前导 \hspace 由 4em 改 3.5em。
%      实测原状目录块在版框内偏下（天头留白 139.6pt / 地脚留白 110.7pt），
%      上移 0.5em = 13.6pt 后上下基本均分，视觉居中。
\NewDocumentCommand { \juan } { mm } {
	\clearpage
	\tl_set:Nn \l__zhvt_chapter_title_tl { #1 }
	\tl_set:Nn \l_tmpa_tl { #2 }
	\tl_set_eq:NN \l_tmpc_tl \l_tmpa_tl
	\int_step_inline:nn { 7 - \str_count:n { #2 } }
	  { \tl_put_right:Nn \l_tmpc_tl { 　 } }
	\tl_set:Nx \l_tmpb_tl { \zhnum { page } }
	\tl_put_left:Nn \l_tmpb_tl { \hspace { 4em } }
	\tl_put_left:Nn \l_tmpb_tl { \l_tmpc_tl }
	\tl_put_left:Nn \l_tmpb_tl { \par \hspace { 3.5em } }
	\addtocontents { toc } { \l_tmpb_tl }
	\skip_horizontal:n { \l__zhvt_font_size_dim * 4 }
	\l_tmpa_tl
	\par
	\thispagestyle { empty }
}
\cs_set_protected:Npn \__zhvt_set_shukou_content: {
	% ★ 鱼尾上方留白（title 静态框）嵌入闲叟钤印：
	%   整页 /Rotate 90，故此处把图用 \rotatebox{90} 转正为竖章；
	%   尺寸卡在版心条幅内（书内约 7mm 宽 × 33.8mm 高，留余量不压鱼尾）。
	\setstaticcontents * { title } {
		\vbox_to_ht:nn { \l__zhvt_baseline_skip_dim } {
			\vskip 14.2pt
			\hbox_to_wd:nn { \l__zhvt_shukou_title_height_dim } {
				\hfil
				\rotatebox [ origin = c ] { 90 } { \includegraphics [ width = 7mm , height = 33.8mm ] { yshs.png } }
				\hfil
			}
		}
	}
% ★ 版心（书口）统一改法（2026-11-16 愚大人定）：18pt＋框内居中＋书口文字同版框色
%   ① 原 27.2pt 时 7 字长篇名（如「廉颇蔺相如列传」）末字溢出 chapter 框、压住页码；
%      改 \fontsize{18pt}{18pt}，前后 \hfil 令字串在框内居中（不随篇名字数漂移）。
%   ② \vbox{\vskip 8.2pt} 定量下移，抵消缩字号后栏宽方向 8.2pt 上偏（像素实测 y 中心 543.3 = 原版）。
%      ⚠ \raisebox / \vspace* 在 flowfram 框内无效；换字号须同步换补偿：16pt→10.1、19pt→7.4。
%   ③ chapter 与 page~number 均置 \color{sealbrown}（HTML 874434），与版框／界行同色；正文仍黑。
%   ⚠ 页码字位、钤印、鱼尾、界行、正文一律未动；五部合订本同此口径。
	\setdynamiccontents * { chapter }
	{ \color{sealbrown} \vbox { \vskip 8.2pt \hbox to \l__zhvt_shukou_chapter_title_height_dim
		{ \hfil \fontsize{18pt}{18pt}\selectfont \l__zhvt_title_str \quad \l__zhvt_chapter_title_tl \hfil } } }
	\setdynamiccontents * { page~number } { \color{sealbrown} \quad \zhnum { page } }
}
\__zhvt_set_shukou_content:
\ExplSyntaxOff

\maketitle

% 单篇，无多卷目录，不渲染目录页

\juan{楚世家}{卷四十~楚世家第十}
\hspace*{2em}楚之先祖出自帝颛顼高阳。高阳者。黄帝之孙。昌意之子也。高阳生称。称生卷章。卷章生重黎。重黎为帝喾高辛居火正。甚有功。能光融天下。帝喾命曰祝融。共工氏作乱。帝喾使重黎诛之而不尽。帝乃以庚寅日诛重黎。而以其弟吴回为重黎后。复居火正。为祝融。\\
\hspace*{2em}吴回生陆终。陆终生子六人。坼剖而产焉。其长一曰昆吾。二曰参胡。三曰彭祖。四曰会人。五曰曹姓。六曰季连。芈姓。楚其后也。昆吾氏。夏之时尝为侯伯。桀之时汤灭之。彭祖氏。殷之时尝为侯伯。殷之末世灭彭祖氏。季连生附沮。附沮生穴熊。其后中微。或在中国。或在蛮夷。弗能纪其世。\\
\hspace*{2em}周文王之时。季连之苗裔曰鬻熊。鬻熊子事文王。蚤卒。其子曰熊丽。熊丽生熊狂。熊狂生熊绎。\\
\hspace*{2em}熊绎当周成王之时。举文。武勤劳之后嗣。而封熊绎于楚蛮。封以子男之田。姓芈氏。居丹阳。楚子熊绎与鲁公伯禽。卫康叔子牟。晋侯燮。齐太公子吕伋俱事成王。\\
\hspace*{2em}熊绎生熊艾。熊艾生熊䵣。熊䵣生熊胜。熊胜以弟熊杨为后。熊杨生熊渠。\\
\hspace*{2em}熊渠生子三人。当周夷王之时。王室微。诸侯或不朝。相伐。熊渠甚得江汉间民和。乃兴兵伐庸。杨粤。至于鄂。熊渠曰。﹁我蛮夷也。不与中国之号谥。﹂乃立其长子康为句亶王。中子红为鄂王。少子执疵为越章王。皆在江上楚蛮之地。及周厉王之时。暴虐。熊渠畏其伐楚。亦去其王。\\
\hspace*{2em}后为熊毋康。毋康蚤死。熊渠卒。子熊挚红立。挚红卒。其弟弑而代立。曰熊延。熊延生熊勇。\\
\hspace*{2em}熊勇六年。而周人作乱。攻厉王。厉王出奔彘。熊勇十年。卒。弟熊严为后。\\
\hspace*{2em}熊严十年。卒。有子四人。长子伯霜。中子仲雪。次子叔堪。少子季徇。熊严卒。长子伯霜代立。是为熊霜。\\
\hspace*{2em}熊霜元年。周宣王初立。熊霜六年。卒。三弟争立。仲雪死。叔堪亡。避难于濮。而少弟季徇立。是为熊徇。熊徇十六年。郑桓公初封于郑。二十二年。熊徇卒。子熊咢立。熊咢九年。卒。子熊仪立。是为若敖。\\
\hspace*{2em}若敖二十年。周幽王为犬戎所弑。周东徙。而秦襄公始列为诸侯。\\
\hspace*{2em}二十七年。若敖卒。子熊坎立。是为霄敖。霄敖六年。卒。子熊眴立。是为鼢冒。鼢冒十三年。晋始乱。以曲沃之故。鼢冒十七年。卒。鼢冒弟熊通弑鼢冒子而代立。是为楚武王。\\
\hspace*{2em}武王十七年。晋之曲沃庄伯弑主国晋孝侯。十九年。郑伯弟段作乱。二十一年。郑侵天子之田。二十三年。卫弑其君桓公。二十九年。鲁弑其君隐公。三十一年。宋太宰华督弑其君殇公。\\
\hspace*{2em}三十五年。楚伐随。随曰。﹁我无罪。﹂楚曰。﹁我蛮夷也。今诸侯皆为叛相侵。或相杀。我有敝甲。欲以观中国之政。请王室尊吾号。﹂随人为之周。请尊楚。王室不听。还报楚。三十七年。楚熊通怒曰。﹁吾先鬻熊。文王之师也。蚤终。成王举我先公。乃以子男田令居楚。蛮夷皆率服。而王不加位。我自尊耳。﹂乃自立为武王。与随人盟而去。于是始开濮地而有之。\\
\hspace*{2em}五十一年。周召随侯。数以立楚为王。楚怒。以随背己。伐随。武王卒师中而兵罢。子文王熊赀立。始都郢。\\
\hspace*{2em}文王二年。伐申过邓。邓人曰﹁楚王易取﹂。邓侯不许也。六年。伐蔡。虏蔡哀侯以归。已而释之。楚强。陵江汉间小国。小国皆畏之。十一年。齐桓公始霸。楚亦始大。\\
\hspace*{2em}十二年。伐邓。灭之。十三年。卒。子熊囏立。是为庄敖。庄敖五年。欲杀其弟熊恽。恽奔随。与随袭弑庄敖代立。是为成王。\\
\hspace*{2em}成王恽元年。初即位。布德施惠。结旧好于诸侯。使人献天子。天子赐胙。曰。﹁镇尔南方夷越之乱。无侵中国。﹂于是楚地千里。\\
\hspace*{2em}十六年。齐桓公以兵侵楚。至陉山。楚成王使将军屈完以兵御之。与桓公盟。桓公数以周之赋不入王室。楚许之。乃去。\\
\hspace*{2em}十八年。成王以兵北伐许。许君肉袒谢。乃释之。二十二年。伐黄。二十六年。灭英。\\
\hspace*{2em}三十三年。宋襄公欲为盟会。召楚。楚王怒曰。﹁召我。我将好往袭辱之。﹂遂行。至盂。遂执辱宋公。已而归之。三十四年。郑文公南朝楚。楚成王北伐宋。败之泓。射伤宋襄公。襄公遂病创死。\\
\hspace*{2em}三十五年。晋公子重耳过楚。成王以诸侯客礼飨。而厚送之于秦。\\
\hspace*{2em}三十九年。鲁僖公来请兵以伐齐。楚使申侯将兵伐齐。取榖。置齐桓公子雍焉。齐桓公七子皆奔楚。楚尽以为上大夫。灭夔。夔不祀祝融。鬻熊故也。\\
\hspace*{2em}夏。伐宋。宋告急于晋。晋救宋。成王罢归。将军子玉请战。成王曰。﹁重耳亡居外久。卒得反国。天之所开。不可当。﹂子玉固请。乃与之少师而去。晋果败子玉于城濮。成王怒。诛子玉。\\
\hspace*{2em}四十六年。初。成王将以商臣为太子。语令尹子上。子上曰。﹁君之齿未也。而又多内宠。绌乃乱也。楚国之举常在少者。且商臣蜂目而豺声。忍人也。不可立也。﹂王不听。立之。后又欲立子职而绌太子商臣。商臣闻而未审也。告其傅潘崇曰。﹁何以得其实。﹂崇曰。﹁飨王之宠姬江芈而勿敬也。﹂商臣从之。江芈怒曰。﹁宜乎王之欲杀若而立职也。﹂商臣告潘崇曰。﹁信矣。﹂崇曰。﹁能事之乎。﹂曰。﹁不能。﹂﹁能亡去乎。﹂曰。﹁不能。﹂﹁能行大事乎。﹂曰。﹁能。﹂冬十月。商臣以宫卫兵围成王。成王请食熊蹯而死。不听。丁未。成王自绞杀。商臣代立。是为穆王。\\
\hspace*{2em}穆王立。以其太子宫予潘崇。使为太师。掌国事。穆王三年。灭江。四年。灭六。蓼。六。蓼。皋陶之后。八年。伐陈。十二年。卒。子庄王侣立。\\
\hspace*{2em}庄王即位三年。不出号令。日夜为乐。令国中曰。﹁有敢谏者死无赦。﹂伍举入谏。庄王左抱郑姬。右抱越女。坐钟鼓之间。伍举曰。﹁愿有进隐。﹂曰。﹁有鸟在于阜。三年不蜚不鸣。是何鸟也。﹂庄王曰。﹁三年不蜚。蜚将冲天。三年不鸣。鸣将惊人。举退矣。吾知之矣。﹂居数月。淫益甚。大夫苏从乃入谏。王曰。﹁若不闻令乎。﹂对曰。﹁杀身以明君。臣之愿也。﹂于是乃罢淫乐。听政。所诛者数百人。所进者数百人。任伍举。苏从以政。国人大说。是岁灭庸。六年。伐宋。获五百乘。\\
\hspace*{2em}八年。伐陆浑戎。遂至洛。观兵于周郊。周定王使王孙满劳楚王。楚王问鼎小大轻重。对曰。﹁在德不在鼎。﹂庄王曰。﹁子无阻九鼎。楚国折钩之喙。足以为九鼎。﹂王孙满曰。﹁呜呼。君王其忘之乎。昔虞夏之盛。远方皆至。贡金九牧。铸鼎象物。百物而为之备。使民知神奸。桀有乱德。鼎迁于殷。载祀六百。殷纣暴虐。鼎迁于周。德之休明。虽小必重。其奸回昏乱。虽大必轻。昔成王定鼎于郏鄏。卜世三十。卜年七百。天所命也。周德虽衰。天命未改。鼎之轻重。未可问也。﹂楚王乃归。\\
\hspace*{2em}九年。相若敖氏。人或谗之王。恐诛。反攻王。王击灭若敖氏之族。十三年。灭舒。\\
\hspace*{2em}十六年。伐陈。杀夏征舒。征舒弑其君。故诛之也。已破陈。即县之。群臣皆贺。申叔时使齐来。不贺。王问。对曰。﹁鄙语曰。牵牛径人田。田主取其牛。径者则不直矣。取之牛不亦甚乎。且王以陈之乱而率诸侯伐之。以义伐之而贪其县。亦何以复令于天下。﹂庄王乃复国陈后。\\
\hspace*{2em}十七年春。楚庄王围郑。三月克之。入自皇门。郑伯肉袒牵羊以逆。曰。﹁孤不天。不能事君。君用怀怒。以及敝邑。孤之罪也。敢不惟命是听。宾之南海。若以臣妾赐诸侯。亦惟命是听。若君不忘厉。宣。桓。武。不绝其社稷。使改事君。孤之愿也。非所敢望也。敢布腹心。﹂楚群臣曰。﹁王勿许。﹂庄王曰。﹁其君能下人。必能信用其民。庸可绝乎。﹂庄王自手旗。左右麾军。引兵去三十里而舍。遂许之平。潘尪入盟。子良出质。夏六月。晋救郑。与楚战。大败晋师河上。遂至衡雍而归。\\
\hspace*{2em}二十年。围宋。以杀楚使也。围宋五月。城中食尽。易子而食。析骨而炊。宋华元出告以情。庄王曰。﹁君子哉。﹂遂罢兵去。\\
\hspace*{2em}二十三年。庄王卒。子共王审立。\\
\hspace*{2em}共王十六年。晋伐郑。郑告急。共王救郑。与晋兵战鄢陵。晋败楚。射中共王目。共王召将军子反。子反嗜酒。从者竖阳榖进酒。醉。王怒。射杀子反。遂罢兵归。\\
\hspace*{2em}三十一年。共王卒。子康王招立。康王立十五年卒。子员立。是为郏敖。\\
\hspace*{2em}康王宠弟公子围。子比。子皙。弃疾。郏敖三年。以其季父康王弟公子围为令尹。主兵事。四年。围使郑。道闻王疾而还。十二月己酉。围入问王疾。绞而弑之。遂杀其子莫及平夏。使使赴于郑。伍举问曰。﹁谁为后。﹂对曰。﹁寡大夫围。﹂伍举更曰。﹁共王之子围为长。﹂子比奔晋。而围立。是为灵王。\\
\hspace*{2em}灵王三年六月。楚使使告晋。欲会诸侯。诸侯皆会楚于申。伍举曰。﹁昔夏启有钧台之飨。商汤有景亳之命。周武王有盟津之誓。成王有岐阳之蒐。康王有丰宫之朝。穆王有涂山之会。齐桓有召陵之师。晋文有践土之盟。君其何用。﹂灵王曰。﹁用桓公。﹂时郑子产在焉。于是晋。宋。鲁。卫不往。灵王已盟。有骄色。伍举曰。﹁桀为有仍之会。有缗叛之。纣为黎山之会。东夷叛之。幽王为太室之盟。戎。翟叛之。君其慎终。﹂\\
\hspace*{2em}七月。楚以诸侯兵伐吴。围朱方。八月。克之。囚庆封。灭其族。以封徇。曰。﹁无效齐庆封弑其君而弱其孤。以盟诸大夫。﹂封反曰。﹁莫如楚共王庶子围弑其君兄之子员而代之立。﹂于是灵王使疾杀之。\\
\hspace*{2em}七年。就章华台。下令内亡人实之。\\
\hspace*{2em}八年。使公子弃疾将兵灭陈。十年。召蔡侯。醉而杀之。使弃疾定蔡。因为陈蔡公。\\
\hspace*{2em}十一年。伐徐以恐吴。灵王次于干溪以待之。王曰。﹁齐。晋。鲁。卫。其封皆受宝器。我独不。今吾使使周求鼎以为分。其予我乎。﹂析父对曰。﹁其予君王哉。昔我先王熊绎辟在荆山。荜露蓝蒌。以处草莽。跋涉山林以事天子。唯是桃弧棘矢以共王事。齐。王舅也。晋及鲁。卫。王母弟也。楚是以无分而彼皆有。周今与四国服事君王。将惟命是从。岂敢爱鼎。﹂灵王曰。﹁昔我皇祖伯父昆吾旧许是宅。今郑人贪其田。不我予。今我求之。其予我乎。﹂对曰。﹁周不爱鼎。郑安敢爱田。﹂灵王曰。﹁昔诸侯远我而畏晋。今吾大城陈。蔡。不羹。赋皆千乘。诸侯畏我乎。﹂对曰。﹁畏哉。﹂灵王喜曰。﹁析父善言古事焉。﹂\\
\hspace*{2em}十二年春。楚灵王乐干溪。不能去也。国人苦役。初。灵王会兵于申。僇越大夫常寿过。杀蔡大夫观起。起子从亡在吴。乃劝吴王伐楚。为间越大夫常寿过而作乱。为吴间。使矫公子弃疾命召公子比于晋。至蔡。与吴。越兵欲袭蔡。令公子比见弃疾。与盟于邓。遂入杀灵王太子禄。立子比为王。公子子皙为令尹。弃疾为司马。先除王宫。观从从师于干溪。令楚众曰。﹁国有王矣。先归。复爵邑田室。后者迁之。﹂楚众皆溃。去灵王而归。\\
\hspace*{2em}灵王闻太子禄之死也。自投车下。而曰。﹁人之爱子亦如是乎。﹂侍者曰。﹁甚是。﹂王曰。﹁余杀人之子多矣。能无及此乎。﹂右尹曰。﹁请待于郊以听国人。﹂王曰。﹁众怒不可犯。﹂曰。﹁且入大县而乞师于诸侯。﹂王曰。﹁皆叛矣。﹂又曰。﹁且奔诸侯以听大国之虑。﹂王曰。﹁大福不再。祗取辱耳。﹂于是王乘舟将欲入鄢。右尹度王不用其计。惧俱死。亦去王亡。\\
\hspace*{2em}灵王于是独傍徨山中。野人莫敢入王。王行遇其故鋗人。谓曰。﹁为我求食。我已不食三日矣。﹂鋗人曰。﹁新王下法。有敢饷王从王者。罪及三族。且又无所得食。﹂王因枕其股而卧。鋗人又以土自代。逃去。王觉而弗见。遂饥弗能起。芋尹申无宇之子申亥曰。﹁吾父再犯王命。王弗诛。恩孰大焉。﹂乃求王。遇王饥于厘泽。奉之以归。夏五月癸丑。王死申亥家。申亥以二女从死。并葬之。\\
\hspace*{2em}是时楚国虽已立比为王。畏灵王复来。又不闻灵王死。故观从谓初王比曰。﹁不杀弃疾。虽得国犹受祸。﹂王曰。﹁余不忍。﹂从曰。﹁人将忍王。﹂王不听。乃去。弃疾归。国人每夜惊。曰。﹁灵王入矣。﹂乙卯夜。弃疾使船人从江上走呼曰。﹁灵王至矣。﹂国人愈惊。又使曼成然告初王比及令尹子皙曰。﹁王至矣。国人将杀君。司马将至矣。君蚤自图。无取辱焉。众怒如水火。不可救也。﹂初王及子皙遂自杀。丙辰。弃疾即位为王。改名熊居。是为平王。\\
\hspace*{2em}平王以诈弑两王而自立。恐国人及诸侯叛之。乃施惠百姓。复陈蔡之地而立其后如故。归郑之侵地。存恤国中。修政教。吴以楚乱故。获五率以归。平王谓观从。﹁恣尔所欲。﹂欲为卜尹。王许之。\\
\hspace*{2em}初。共王有宠子五人。无适立。乃望祭群神。请神决之。使主社稷。而阴与巴姬埋璧于室内。召五公子斋而入。康王跨之。灵王肘加之。子比。子皙皆远之。平王幼。抱其上而拜。压纽。故康王以长立。至其子失之。围为灵王。及身而弑。子比为王十余日。子皙不得立。又俱诛。四子皆绝无后。唯独弃疾后立。为平王。竟续楚祀。如其神符。\\
\hspace*{2em}初。子比自晋归。韩宣子问叔向曰。﹁子比其济乎。﹂对曰。﹁不就。﹂宣子曰。﹁同恶相求。如市贾焉。何为不就。﹂对曰。﹁无与同好。谁与同恶。取国有五难。有宠无人。一也。有人无主。二也。有主无谋。三也。有谋而无民。四也。有民而无德。五也。子比在晋十三年矣。晋。楚之从不闻通者。可谓无人矣。族尽亲叛。可谓无主矣。无衅而动。可谓无谋矣。为羁终世。可谓无民矣。亡无爱征。可谓无德矣。王虐而不忌。子比涉五难以弑君。谁能济之。有楚国者。其弃疾乎。君陈。蔡。方城外属焉。苛慝不作。盗贼伏隐。私欲不违。民无怨心。先神命之。国民信之。芈姓有乱。必季实立。楚之常也。子比之官。则右尹也。数其贵宠。则庶子也。以神所命。则又远之。民无怀焉。将何以立。﹂宣子曰。﹁齐桓。晋文不亦是乎。﹂对曰。﹁齐桓。卫姬之子也。有宠于厘公。有鲍叔牙。宾须无。隰朋以为辅。有莒。卫以为外主。有高。国以为内主。从善如流。施惠不倦。有国。不亦宜乎。昔我文公。狐季姬之子也。有宠于献公。好学不倦。生十七年。有士五人。有先大夫子余。子犯以为腹心。有魏犨。贾佗以为股肱。有齐。宋。秦。楚以为外主。有栾。郄。狐。先以为内主。亡十九年。守志弥笃。惠。怀弃民。民从而与之。故文公有国。不亦宜乎。子比无施于民。无援于外。去晋。晋不送。归楚。楚不迎。何以有国。﹂子比果不终焉。卒立者弃疾。如叔向言也。\\
\hspace*{2em}平王二年。使费无忌如秦为太子建取妇。妇好。来。未至。无忌先归。说平王曰。﹁秦女好。可自娶。为太子更求。﹂平王听之。卒自娶秦女。生熊珍。更为太子娶。是时伍奢为太子太傅。无忌为少傅。无忌无宠于太子。常谗恶太子建。建时年十五矣。其母蔡女也。无宠于王。王稍益疏外建也。\\
\hspace*{2em}六年。使太子建居城父。守边。无忌又日夜谗太子建于王曰。﹁自无忌入秦女。太子怨。亦不能无望于王。王少自备焉。且太子居城父。擅兵。外交诸侯。且欲入矣。﹂平王召其傅伍奢责之。伍奢知无忌谗。乃曰。﹁王奈何以小臣疏骨肉。﹂无忌曰。﹁今不制。后悔也。﹂于是王遂囚伍奢。乃令司马奋扬召太子建。欲诛之。太子闻之。亡奔宋。\\
\hspace*{2em}无忌曰。﹁伍奢有二子。不杀者为楚国患。盍以免其父召之。必至。﹂于是王使使谓奢。﹁能致二子则生。不能将死。﹂奢曰。﹁尚至。胥不至。﹂王曰。﹁何也。﹂奢曰。﹁尚之为人。廉。死节。慈孝而仁。闻召而免父。必至。不顾其死。胥之为人。智而好谋。勇而矜功。知来必死。必不来。然为楚国忧者必此子。﹂于是王使人召之。曰。﹁来。吾免尔父。﹂伍尚谓伍胥曰。﹁闻父免而莫奔。不孝也。父戮莫报。无谋也。度能任事。知也。子其行矣。我其归死。﹂伍尚遂归。伍胥弯弓属矢。出见使者。曰。﹁父有罪。何以召其子为。﹂将射。使者还走。遂出奔吴。伍奢闻之。曰。﹁胥亡。楚国危哉。﹂楚人遂杀伍奢及尚。\\
\hspace*{2em}十年。楚太子建母在居巢。开吴。吴使公子光伐楚。遂败陈。蔡。取太子建母而去。楚恐。城郢。初。吴之边邑卑梁与楚边邑钟离小童争桑。两家交怒相攻。灭卑梁人。卑梁大夫怒。发邑兵攻钟离。楚王闻之怒。发国兵灭卑梁。吴王闻之大怒。亦发兵。使公子光因建母家攻楚。遂灭钟离。居巢。楚乃恐而城郢。\\
\hspace*{2em}十三年。平王卒。将军子常曰。﹁太子珍少。且其母乃前太子建所当娶也。﹂欲立令尹子西。子西。平王之庶弟也。有义。子西曰。﹁国有常法。更立则乱。言之则致诛。﹂乃立太子珍。是为昭王。\\
\hspace*{2em}昭王元年。楚众不说费无忌。以其谗亡太子建。杀伍奢子父与郄宛。宛之宗姓伯氏子嚭及子胥皆奔吴。吴兵数侵楚。楚人怨无忌甚。楚令尹子常诛无忌以说众。众乃喜。\\
\hspace*{2em}四年。吴三公子奔楚。楚封之以扞吴。五年。吴伐取楚之六。潜。七年。楚使子常伐吴。吴大败楚于豫章。\\
\hspace*{2em}十年冬。吴王阖闾。伍子胥。伯嚭与唐。蔡俱伐楚。楚大败。吴兵遂入郢。辱平王之墓。以伍子胥故也。吴兵之来。楚使子常以兵迎之。夹汉水阵。吴伐败子常。子常亡奔郑。楚兵走。吴乘胜逐之。五战及郢。己卯。昭王出奔。庚辰。吴人入郢。\\
\hspace*{2em}昭王亡也至云梦。云梦不知其王也。射伤王。王走郧。郧公之弟怀曰。﹁平王杀吾父。今我杀其子。不亦可乎。﹂郧公止之。然恐其弑昭王。乃与王出奔随。吴王闻昭王往。即进击随。谓随人曰。﹁周之子孙封于江汉之间者。楚尽灭之。﹂欲杀昭王。王从臣子綦乃深匿王。自以为王。谓随人曰。﹁以我予吴。﹂随人卜予吴。不吉。乃谢吴王曰。﹁昭王亡。不在随。﹂吴请入自索之。随不听。吴亦罢去。\\
\hspace*{2em}昭王之出郢也。使申鲍胥请救于秦。秦以车五百乘救楚。楚亦收余散兵。与秦击吴。十一年六月。败吴于稷。会吴王弟夫概见吴王兵伤败。乃亡归。自立为王。阖闾闻之。引兵去楚。归击夫概。夫概败。奔楚。楚封之堂谿。号为堂谿氏。\\
\hspace*{2em}楚昭王灭唐。九月。归入郢。十二年。吴复伐楚。取番。楚恐。去郢。北徙都鄀。\\
\hspace*{2em}十六年。孔子相鲁。二十年。楚灭顿。灭胡。二十一年。吴王阖闾伐越。越王句践射伤吴王。遂死。吴由此怨越而不西伐楚。\\
\hspace*{2em}二十七年春。吴伐陈。楚昭王救之。军城父。十月。昭王病于军中。有赤云如鸟。夹日而蜚。昭王问周太史。太史曰。﹁是害于楚王。然可移于将相。﹂将相闻是言。乃请自以身祷于神。昭王曰。﹁将相。孤之股肱也。今移祸。庸去是身乎。﹂弗听。卜而河为祟。大夫请祷河。昭王曰。﹁自吾先王受封。望不过江。汉。而河非所获罪也。﹂止不许。孔子在陈。闻是言。曰。﹁楚昭王通大道矣。其不失国。宜哉。﹂\\
\hspace*{2em}昭王病甚。乃召诸公子大夫曰。﹁孤不佞。再辱楚国之师。今乃得以天寿终。孤之幸也。﹂让其弟公子申为王。不可。又让次弟公子结。亦不可。乃又让次弟公子闾。五让。乃后许为王。将战。庚寅。昭王卒于军中。子闾曰。﹁王病甚。舍其子让群臣。臣所以许王。以广王意也。今君王卒。臣岂敢忘君王之意乎。﹂乃与子西。子綦谋。伏师闭涂。迎越女之子章立之。是为惠王。然后罢兵归。葬昭王。\\
\hspace*{2em}惠王二年。子西召故平王太子建之子胜于吴。以为巢大夫。号曰白公。白公好兵而下士。欲报仇。六年。白公请兵令尹子西伐郑。初。白公父建亡在郑。郑杀之。白公亡走吴。子西复召之。故以此怨郑。欲伐之。子西许而未为发兵。八年。晋伐郑。郑告急楚。楚使子西救郑。受赂而去。白公胜怒。乃遂与勇力死士石乞等袭杀令尹子西。子綦于朝。因劫惠王。置之高府。欲弑之。惠王从者屈固负王亡走昭王夫人宫。白公自立为王。月余。会叶公来救楚。楚惠王之徒与共攻白公。杀之。惠王乃复位。是岁也。灭陈而县之。\\
\hspace*{2em}十三年。吴王夫差强。陵齐。晋。来伐楚。十六年。越灭吴。四十二年。楚灭蔡。四十四年。楚灭杞。与秦平。是时越已灭吴而不能正江。淮北。楚东侵。广地至泗上。\\
\hspace*{2em}五十七年。惠王卒。子简王中立。\\
\hspace*{2em}简王元年。北伐灭莒。八年。魏文侯。韩武子。赵桓子始列为诸侯。\\
\hspace*{2em}二十四年。简王卒。子声王当立。声王六年。盗杀声王。子悼王熊疑立。悼王二年。三晋来伐楚。至乘丘而还。四年。楚伐周。郑杀子阳。九年。伐韩。取负黍。十一年。三晋伐楚。败我大梁。榆关。楚厚赂秦。与之平。二十一年。悼王卒。子肃王臧立。\\
\hspace*{2em}肃王四年。蜀伐楚。取兹方。于是楚为扞关以距之。十年。魏取我鲁阳。十一年。肃王卒。无子。立其弟熊良夫。是为宣王。\\
\hspace*{2em}宣王六年。周天子贺秦献公。秦始复强。而三晋益大。魏惠王。齐威王尤强。三十年。秦封卫鞅于商。南侵楚。是年。宣王卒。子威王熊商立。\\
\hspace*{2em}威王六年。周显王致文武胙于秦惠王。\\
\hspace*{2em}七年。齐孟尝君父田婴欺楚。楚威王伐齐。败之于徐州。而令齐必逐田婴。田婴恐。张丑伪谓楚王曰。﹁王所以战胜于徐州者。田盼子不用也。盼子者。有功于国。而百姓为之用。婴子弗善而用申纪。申纪者。大臣不附。百姓不为用。故王胜之也。今王逐婴子。婴子逐。盼子必用矣。复搏其士卒以与王遇。必不便于王矣。﹂楚王因弗逐也。\\
\hspace*{2em}十一年。威王卒。子怀王熊槐立。魏闻楚丧。伐楚。取我陉山。\\
\hspace*{2em}怀王元年。张仪始相秦惠王。四年。秦惠王初称王。\\
\hspace*{2em}六年。楚使柱国昭阳将兵而攻魏。破之于襄陵。得八邑。又移兵而攻齐。齐王患之。陈轸适为秦使齐。齐王曰。﹁为之奈何。﹂陈轸曰。﹁王勿忧。请令罢之。﹂即往见昭阳军中。曰。﹁愿闻楚国之法。破军杀将者何以贵之。﹂昭阳曰。﹁其官为上柱国。封上爵执圭。﹂陈轸曰。﹁其有贵于此者乎。﹂昭阳曰。﹁令尹。﹂陈轸曰。﹁今君已为令尹矣。此国冠之上。臣请得譬之。人有遗其舍人一卮酒者。舍人相谓曰。﹃数人饮此。不足以遍。请遂画地为蛇。蛇先成者独饮之。﹄一人曰。﹃吾蛇先成。﹄举酒而起。曰。﹃吾能为之足。﹄及其为之足。而后成人夺之酒而饮之。曰。﹃蛇固无足。今为之足。是非蛇也。﹄今君相楚而攻魏。破军杀将。功莫大焉。冠之上不可以加矣。今又移兵而攻齐。攻齐胜之。官爵不加于此。攻之不胜。身死爵夺。有毁于楚。此为蛇为足之说也。不若引兵而去以德齐。此持满之术也。﹂昭阳曰。﹁善。﹂引兵而去。\\
\hspace*{2em}燕。韩君初称王。秦使张仪与楚。齐。魏相会。盟啮桑。\\
\hspace*{2em}十一年。苏秦约从山东六国共攻秦。楚怀王为从长。至函谷关。秦出兵击六国。六国兵皆引而归。齐独后。十二年。齐愍王伐败赵。魏军。秦亦伐败韩。与齐争长。\\
\hspace*{2em}十六年。秦欲伐齐。而楚与齐从亲。秦惠王患之。乃宣言张仪免相。使张仪南见楚王。谓楚王曰。﹁敝邑之王所甚说者无先大王。虽仪之所甚愿为门阑之厮者亦无先大王。敝邑之王所甚憎者无先齐王。虽仪之所甚憎者亦无先齐王。而大王和之。是以敝邑之王不得事王。而令仪亦不得为门阑之厮也。王为仪闭关而绝齐。今使使者从仪西取故秦所分楚商于之地方六百里。如是则齐弱矣。是北弱齐。西德于秦。私商于以为富。此一计而三利俱至也。﹂怀王大悦。乃置相玺于张仪。日与置酒。宣言﹁吾复得吾商于之地﹂。群臣皆贺。而陈轸独吊。怀王曰。﹁何故。﹂陈轸对曰。﹁秦之所为重王者。以王之有齐也。今地未可得而齐交先绝。是楚孤也。夫秦又何重孤国哉。必轻楚矣。且先出地而后绝齐。则秦计不为。先绝齐而后责地。则必见欺于张仪。见欺于张仪。则王必怨之。怨之。是西起秦患。北绝齐交。西起秦患。北绝齐交。则两国之兵必至。臣故吊。﹂楚王弗听。因使一将军西受封地。\\
\hspace*{2em}张仪至秦。详醉坠车。称病不出三月。地不可得。楚王曰。﹁仪以吾绝齐为尚薄邪。﹂乃使勇士宋遗北辱齐王。齐王大怒。折楚符而合于秦。秦齐交合。张仪乃起朝。谓楚将军曰。﹁子何不受地。从某至某。广袤六里。﹂楚将军曰。﹁臣之所以见命者六百里。不闻六里。﹂即以归报怀王。怀王大怒。兴师将伐秦。陈轸又曰。﹁伐秦非计也。不如因赂之一名都。与之伐齐。是我亡于秦。取偿于齐也。吾国尚可全。今王已绝于齐而责欺于秦。是吾合秦齐之交而来天下之兵也。国必大伤矣。﹂楚王不听。遂绝和于秦。发兵西攻秦。秦亦发兵击之。\\
\hspace*{2em}十七年春。与秦战丹阳。秦大败我军。斩甲士八万。虏我大将军屈匄。裨将军逢侯丑等七十余人。遂取汉中之郡。楚怀王大怒。乃悉国兵复袭秦。战于蓝田。大败楚军。韩。魏闻楚之困。乃南袭楚。至于邓。楚闻。乃引兵归。\\
\hspace*{2em}十八年。秦使使约复与楚亲。分汉中之半以和楚。楚王曰。﹁愿得张仪。不愿得地。﹂张仪闻之。请之楚。秦王曰。﹁楚且甘心于子。奈何。﹂张仪曰。﹁臣善其左右靳尚。靳尚又能得事于楚王幸姬郑袖。袖所言无不从者。且仪以前使负楚以商于之约。今秦楚大战。有恶。臣非面自谢楚不解。且大王在。楚不宜敢取仪。诚杀仪以便国。臣之愿也。﹂仪遂使楚。\\
\hspace*{2em}至。怀王不见。因而囚张仪。欲杀之。仪私于靳尚。靳尚为请怀王曰。﹁拘张仪。秦王必怒。天下见楚无秦。必轻王矣。﹂又谓夫人郑袖曰。﹁秦王甚爱张仪。而王欲杀之。今将以上庸之地六县赂楚。以美人聘楚王。以宫中善歌者为之媵。楚王重地。秦女必贵。而夫人必斥矣。夫人不若言而出之。﹂郑袖卒言张仪于王而出之。仪出。怀王因善遇仪。仪因说楚王以叛从约而与秦合亲。约婚姻。张仪已去。屈原使从齐来。谏王曰。﹁何不诛张仪。﹂怀王悔。使人追仪。弗及。是岁。秦惠王卒。\\
\hspace*{2em}二十年。齐愍王欲为从长。恶楚之与秦合。乃使使遗楚王书曰。﹁寡人患楚之不察于尊名也。今秦惠王死。武王立。张仪走魏。樗里疾。公孙衍用。而楚事秦。夫樗里疾善乎韩。而公孙衍善乎魏。楚必事秦。韩。魏恐。必因二人求合于秦。则燕。赵亦宜事秦。四国争事秦。则楚为郡县矣。王何不与寡人并力收韩。魏。燕。赵。与为从而尊周室。以案兵息民。令于天下。莫敢不乐听。则王名成矣。王率诸侯并伐。破秦必矣。王取武关。蜀。汉之地。私吴。越之富而擅江海之利。韩。魏割上党。西薄函谷。则楚之强百万也。且王欺于张仪。亡地汉中。兵锉蓝田。天下莫不代王怀怒。今乃欲先事秦。愿大王孰计之。﹂\\
\hspace*{2em}楚王业已欲和于秦。见齐王书。犹豫不决。下其议群臣。群臣或言和秦。或曰听齐。昭雎曰。﹁王虽东取地于越。不足以刷耻。必且取地于秦。而后足以刷耻于诸侯。王不如深善齐。韩以重樗里疾。如是则王得韩。齐之重以求地矣。秦破韩宜阳。而韩犹复事秦者。以先王墓在平阳。而秦之武遂去之七十里。以故尤畏秦。不然。秦攻三川。赵攻上党。楚攻河外。韩必亡。楚之救韩。不能使韩不亡。然存韩者楚也。韩已得武遂于秦。以河山为塞。所报德莫如楚厚。臣以为其事王必疾。齐之所信于韩者。以韩公子眛为齐相也。韩已得武遂于秦。王甚善之。使之以齐。韩重樗里疾。疾得齐。韩之重。其主弗敢弃疾也。今又益之以楚之重。樗里子必言秦。复与楚之侵地矣。﹂于是怀王许之。竟不合秦。而合齐以善韩。\\
\hspace*{2em}二十四年。倍齐而合秦。秦昭王初立。乃厚赂于楚。楚往迎妇。二十五年。怀王入与秦昭王盟。约于黄棘。秦复与楚上庸。二十六年。齐。韩。魏为楚负其从亲而合于秦。三国共伐楚。楚使太子入质于秦而请救。秦乃遣客卿通将兵救楚。三国引兵去。\\
\hspace*{2em}二十七年。秦大夫有私与楚太子斗。楚太子杀之而亡归。二十八年。秦乃与齐。韩。魏共攻楚。杀楚将唐眛。取我重丘而去。二十九年。秦复攻楚。大破楚。楚军死者二万。杀我将军景缺。怀王恐。乃使太子为质于齐以求平。三十年。秦复伐楚。取八城。秦昭王遗楚王书曰。﹁始寡人与王约为弟兄。盟于黄棘。太子为质。至欢也。太子陵杀寡人之重臣。不谢而亡去。寡人诚不胜怒。使兵侵君王之边。今闻君王乃令太子质于齐以求平。寡人与楚接境壤界。故为婚姻。所从相亲久矣。而今秦楚不欢。则无以令诸侯。寡人愿与君王会武关。面相约。结盟而去。寡人之愿也。敢以闻下执事。﹂楚怀王见秦王书。患之。欲往。恐见欺。无往。恐秦怒。昭雎曰。﹁王毋行。而发兵自守耳。秦虎狼。不可信。有并诸侯之心。﹂怀王子子兰劝王行。曰。﹁奈何绝秦之欢心。﹂于是往会秦昭王。昭王诈令一将军伏兵武关。号为秦王。楚王至。则闭武关。遂与西至咸阳。朝章台。如蕃臣。不与亢礼。楚怀王大怒。悔不用昭子言。秦因留楚王。要以割巫。黔中之郡。楚王欲盟。秦欲先得地。楚王怒曰。﹁秦诈我而又强要我以地。﹂不复许秦。秦因留之。\\
\hspace*{2em}楚大臣患之。乃相与谋曰。﹁吾王在秦不得还。要以割地。而太子为质于齐。齐。秦合谋。则楚无国矣。﹂乃欲立怀王子在国者。昭雎曰。﹁王与太子俱困于诸侯。而今又倍王命而立其庶子。不宜。﹂乃诈赴于齐。齐愍王谓其相曰。﹁不若留太子以求楚之淮北。﹂相曰。﹁不可。郢中立王。是吾抱空质而行不义于天下也。﹂或曰。﹁不然。郢中立王。因与其新王市曰﹃予我下东国。吾为王杀太子。不然。将与三国共立之﹄。然则东国必可得矣。﹂齐王卒用其相计而归楚太子。太子横至。立为王。是为顷襄王。乃告于秦曰。﹁赖社稷神灵。国有王矣。﹂\\
\hspace*{2em}顷襄王横元年。秦要怀王不可得地。楚立王以应秦。秦昭王怒。发兵出武关攻楚。大败楚军。斩首五万。取析十五城而去。二年。楚怀王亡逃归。秦觉之。遮楚道。怀王恐。乃从间道走赵以求归。赵主父在代。其子惠王初立。行王事。恐。不敢入楚王。楚王欲走魏。秦追至。遂与秦使复之秦。怀王遂发病。顷襄王三年。怀王卒于秦。秦归其丧于楚。楚人皆怜之。如悲亲戚。诸侯由是不直秦。秦楚绝。\\
\hspace*{2em}六年。秦使白起伐韩于伊阙。大胜。斩首二十四万。秦乃遗楚王书曰。﹁楚倍秦。秦且率诸侯伐楚。争一旦之命。愿王之饬士卒。得一乐战。﹂楚顷襄王患之。乃谋复与秦平。七年。楚迎妇于秦。秦楚复平。\\
\hspace*{2em}十一年。齐秦各自称为帝。月余。复归帝为王。\\
\hspace*{2em}十四年。楚顷襄王与秦昭王好会于宛。结和亲。十五年。楚王与秦。三晋。燕共伐齐。取淮北。十六年。与秦昭王好会于鄢。其秋。复与秦王会穰。\\
\hspace*{2em}十八年。楚人有好以弱弓微缴加归雁之上者。顷襄王闻。召而问之。对曰。﹁小臣之好射鶀雁。罗鸗。小矢之发也。何足为大王道也。且称楚之大。因大王之贤。所弋非直此也。昔者三王以弋道德。五霸以弋战国。故秦。魏。燕。赵者。鶀雁也。齐。鲁。韩。卫者。青首也。驺。费。郯。邳者。罗鸗也。外其余则不足射者。见鸟六双。以王何取。王何不以圣人为弓。以勇士为缴。时张而射之。此六双者。可得而囊载也。其乐非特朝昔之乐也。其获非特凫雁之实也。王朝张弓而射魏之大梁之南。加其右臂而径属之于韩。则中国之路绝而上蔡之郡坏矣。还射圉之东。解魏左肘而外击定陶。则魏之东外弃而大宋。方与二郡者举矣。且魏断二臂。颠越矣。膺击郯国。大梁可得而有也。王綪缴兰台。饮马西河。定魏大梁。此一发之乐也。若王之于弋诚好而不厌。则出宝弓。碆新缴。射噣鸟于东海。还盖长城以为防。朝射东莒。夕发浿丘。夜加即墨。顾据午道。则长城之东收而太山之北举矣。西结境于赵而北达于燕。三国布。则从不待约而可成也。北游目于燕之辽东而南登望于越之会稽。此再发之乐也。若夫泗上十二诸侯。左萦而右拂之。可一旦而尽也。今秦破韩以为长忧。得列城而不敢守也。伐魏而无功。击赵而顾病。则秦魏之勇力屈矣。楚之故地汉中。析。郦可得而复有也。王出宝弓。碆新缴。涉塞。而待秦之倦也。山东。河内可得而一也。劳民休众。南面称王矣。故曰秦为大鸟。负海内而处。东面而立。左臂据赵之西南。右臂傅楚鄢郢。膺击韩魏。垂头中国。处既形便。势有地利。奋翼鼓。方三千里。则秦未可得独招而夜射也。﹂欲以激怒襄王。故对以此言。襄王因召与语。遂言曰。﹁夫先王为秦所欺而客死于外。怨莫大焉。今以匹夫有怨。尚有报万乘。白公。子胥是也。今楚之地方五千里。带甲百万。犹足以踊跃中野也。而坐受困。臣窃为大王弗取也。﹂于是顷襄王遣使于诸侯。复为从。欲以伐秦。秦闻之。发兵来伐楚。\\
\hspace*{2em}楚欲与齐韩连和伐秦。因欲图周。周王赧使武公谓楚相昭子曰。﹁三国以兵割周郊地以便输。而南器以尊楚。臣以为不然。夫弑共主。臣世君。大国不亲。以众胁寡。小国不附。大国不亲。小国不附。不可以致名实。名实不得。不足以伤民。夫有图周之声。非所以为号也。﹂昭子曰。﹁乃图周则无之。虽然。周何故不可图也。﹂对曰。﹁军不五不攻。城不十不围。夫一周为二十晋。公之所知也。韩尝以二十万之众辱于晋之城下。锐士死。中士伤。而晋不拔。公之无百韩以图周。此天下之所知也。夫怨结于两周以塞驺鲁之心。交绝于齐。声失天下。其为事危矣。夫危两周以厚三川。方城之外必为韩弱矣。何以知其然也。西周之地。绝长补短。不过百里。名为天下共主。裂其地不足以肥国。得其众不足以劲兵。虽无攻之。名为弑君。然而好事之君。喜攻之臣。发号用兵。未尝不以周为终始。是何也。见祭器在焉。欲器之至而忘弑君之乱。今韩以器之在楚。臣恐天下以器仇楚也。臣请譬之。夫虎肉臊。其兵利身。人犹攻之也。若使泽中之麋蒙虎之皮。人之攻之必万于虎矣。裂楚之地。足以肥国。诎楚之名。足以尊主。今子将以欲诛残天下之共主。居三代之传器。吞三翮六翼。以高世主。非贪而何。周书曰﹃欲起无先﹄。故器南则兵至矣。﹂于是楚计辍不行。\\
\hspace*{2em}十九年。秦伐楚。楚军败。割上庸。汉北地予秦。二十年。秦将白起拔我西陵。二十一年。秦将白起遂拔我郢。烧先王墓夷陵。楚襄王兵散。遂不复战。东北保于陈城。二十二年。秦复拔我巫。黔中郡。\\
\hspace*{2em}二十三年。襄王乃收东地兵。得十余万。复西取秦所拔我江旁十五邑以为郡。距秦。二十七年。使三万人助三晋伐燕。复与秦平。而入太子为质于秦。楚使左徒侍太子于秦。\\
\hspace*{2em}三十六年。顷襄王病。太子亡归。秋。顷襄王卒。太子熊元代立。是为考烈王。考烈王以左徒为令尹。封以吴。号春申君。\\
\hspace*{2em}考烈王元年。纳州于秦以平。是时楚益弱。\\
\hspace*{2em}六年。秦围邯郸。赵告急楚。楚遣将军景阳救赵。七年。至新中。秦兵去。十二年。秦昭王卒。楚王使春申君吊祠于秦。十六年。秦庄襄王卒。秦王赵政立。二十二年。与诸侯共伐秦。不利而去。楚东徙都寿春。命曰郢。\\
\hspace*{2em}二十五年。考烈王卒。子幽王悍立。李园杀春申君。幽王三年。秦。魏伐楚。秦相吕不韦卒。九年。秦灭韩。十年。幽王卒。同母弟犹代立。是为哀王。哀王立二月余。哀王庶兄负刍之徒袭杀哀王而立负刍为王。是岁。秦虏赵王迁。\\
\hspace*{2em}王负刍元年。燕太子丹使荆轲刺秦王。二年。秦使将军伐楚。大破楚军。亡十余城。三年。秦灭魏。四年。秦将王翦破我军于蕲。而杀将军项燕。\\
\hspace*{2em}五年。秦将王翦。蒙武遂破楚国。虏楚王负刍。灭楚名为郡云。\\
\hspace*{2em}太史公曰。楚灵王方会诸侯于申。诛齐庆封。作章华台。求周九鼎之时。志小天下。及饿死于申亥之家。为天下笑。操行之不得。悲夫。势之于人也。可不慎与。弃疾以乱立。嬖淫秦女。甚乎哉。几再亡国。\\

% ===== 索隐述赞（司马贞《史记索隐》）· 统一样式 =====
% 题头「索隐述赞曰」五字并入反白框；正文每四句一列，定宽盒 \hfil 居列中；
% \fontsize{24}{41} 列距对齐全书界行（121+41k）；\clearpage 使述赞独占一页。
\clearpage
\shuzanhead{索隐述赞曰}\linebreak
{\shuzanbody\fontsize{24}{41}\selectfont
\hbox to \hsize{\hfil 鬻熊之嗣。周封於楚。僻在荆蛮。荜路蓝缕。\hfil}\linebreak
\hbox to \hsize{\hfil 及通而霸。僭号曰武。文既伐申。成亦赦许。\hfil}\linebreak
\hbox to \hsize{\hfil 子圉篡嫡。商臣杀父。天祸未悔。凭奸自怙。\hfil}\linebreak
\hbox to \hsize{\hfil 昭困奔亡。怀迫囚虏。顷襄。考烈。\hfil}\linebreak
\hbox to \hsize{\hfil 祚衰南土。\hfil}%
\par}
\end{document}
